\documentclass[a4paper,11pt,oneside,openany]{book}
\usepackage{exam_style}

\title{CSR101: Python Programming (AI \& Data Engineering Specialization) --- Examination Companion \& Solved Question Papers}
\author{LPU School of Computer Science and Engineering Faculty}
\date{Edition 2026--2027}

\begin{document}

\frontmatter
\begin{titlepage}
\pagecolor{LPUNavy}
\color{white}
\vspace*{1.5cm}

\begin{center}
    {\Huge \bfseries \scshape CSR101: Python Programming}\\[5mm]
    {\LARGE \bfseries AI \& Data Engineering Specialization}\\[3mm]
    {\large \bfseries The Definitive University Examination Companion \& Solved Model Papers}\\[8mm]
    {\color{LPUOrange}\rule{0.8\textwidth}{2pt}}\\[10mm]

    \begin{tcolorbox}[enhanced, width=0.9\textwidth, colback=white, colframe=LPUOrange, boxrule=1.5pt, arc=3mm, halign=center]
        \color{DarkSlate}
        \vspace{2mm}
        {\Large\bfseries Complete Assessment Suite (Sets A \& B for CA-1, Midterm, CA-2, Endterm)}\\[3mm]
        {\large\bfseries Featuring Step-by-Step Code Implementations \& Examiner Marking Rubrics}\\[3mm]
        \begin{tabular}{cc}
            \color{LPUNavy}\textbf{Full CA-1 Papers (Units 1 \& 2)} & \quad \color{LPUNavy}\textbf{Full Midterm Papers (Units 1, 2, 3)} \\
            \color{LPUNavy}\textbf{Full CA-2 Papers (Units 4 \& 5)} & \quad \color{LPUNavy}\textbf{Full End-Term Papers (Units 1 to 6)} \\
        \end{tabular}
        \vspace{2mm}
    \end{tcolorbox}

    \vfill

    {\large \bfseries Prepared for Lovely Professional University (LPU) Computer Science Students}\\[2mm]
    {\small Course Outcomes CO1--CO6 $\bullet$ Bloom's Taxonomy L1--L4 $\bullet$ Exact Marking Rubrics}\\[10mm]

    \begin{tcolorbox}[width=0.65\textwidth, colback=LPUNavy!85!black, colframe=white, boxrule=0.5pt, halign=center]
        \color{white}
        \textbf{LPU Engineering Academic Series}\\[1mm]
        {\footnotesize Academic Session 2026--2027}
    \end{tcolorbox}
    \vspace*{1cm}
\end{center}
\end{titlepage}
\nopagecolor
\newpage

\chapter*{Official Examination Architecture \& Blueprint}
\addcontentsline{toc}{chapter}{Official Examination Architecture \& Blueprint}

This Examination Companion provides verified model question papers and complete step-by-step solutions designed strictly to mirror the Lovely Professional University (LPU) examination structure for \textbf{CSR101: Python Programming (AI \& Data Engineering Specialization)}:

\begin{center}
\begin{tabularx}{\textwidth}{l p{3.2cm} p{2.8cm} c X}
\toprule
\textbf{Assessment} & \textbf{Syllabus Scope} & \textbf{Format} & \textbf{Marks} & \textbf{Instructions \& Rules} \\
\midrule
\textbf{CA-1} & Unit 1 \& Unit 2 & 6 Subjective Qs & 30 & Attempt any 5 out of 6 questions ($5 \times 6\text{M} = 30\text{M}$). Time: 60 mins. \\
\addlinespace
\textbf{Mid-Term} & Unit 1, Unit 2, Unit 3 & 30 MCQs + Subjective & 30 & MCQs with negative marking ($-0.25$ per wrong answer) + short coding problems. Time: 90 mins. \\
\addlinespace
\textbf{CA-2} & Unit 4 \& Unit 5 & 6 Subjective Qs & 30 & Attempt any 5 out of 6 questions ($5 \times 6\text{M} = 30\text{M}$). Time: 60 mins. \\
\addlinespace
\textbf{End-Term} & Units 1 to 6 & \textbf{Part A}: MCQs \newline \textbf{Part B}: Short Coding \newline \textbf{Part C}: Long Projects & 70 & Comprehensive university evaluation spanning syntax, OOP, NumPy, and Pandas pipelines. Time: 3 Hours. \\
\bottomrule
\end{tabularx}
\end{center}

\vspace{5mm}
\noindent
\textbf{Course Outcome Mapping:}
\begin{itemize}[leftmargin=1.5em, itemsep=2pt]
    \item \textbf{CO1 (Unit 1):} Understand CPython internals, memory model, and primitive types.
    \item \textbf{CO2 (Unit 2):} Apply control flow, functions, regular expressions, and algorithmic paradigms.
    \item \textbf{CO3 (Unit 3):} Design high-performance object-oriented structures with decorators and custom dunder methods.
    \item \textbf{CO4 (Unit 4):} Architect modular software systems with custom exception hierarchies and context managers.
    \item \textbf{CO5 (Unit 5):} Construct vectorized matrix algorithms and linear algebra pipelines in NumPy.
    \item \textbf{CO6 (Unit 6):} Implement end-to-end ETL data loading and statistical visualization pipelines.
\end{itemize}
\newpage

\chapter*{Rapid Formula \& Syntax Quick-Reference Sheet}
\addcontentsline{toc}{chapter}{Rapid Formula \& Syntax Quick-Reference Sheet}

\begin{multicols}{2}
\subsection*{Core Syntax \& Precedence}
\textbf{Highest to Lowest:}
\begin{enumerate}[leftmargin=1.2em, itemsep=1pt]
    \item Parentheses, indexing \texttt{x[i]}, call \texttt{f()}
    \item Exponentiation \texttt{**} (Right-to-left!)
    \item Unary \texttt{+x}, \texttt{-x}, \texttt{\textasciitilde x}
    \item \texttt{*}, \texttt{/}, \texttt{//}, \texttt{\%}
    \item \texttt{+}, \texttt{-}
    \item Bitwise shifts \texttt{<<}, \texttt{>>}
    \item Bitwise \texttt{\&}, \texttt{\textasciicircum}, \texttt{\textbar}
    \item Comparisons \texttt{==}, \texttt{!=}, \texttt{<}, \texttt{<=}, \texttt{is}, \texttt{in}
    \item Boolean \texttt{not}, \texttt{and}, \texttt{or}
\end{enumerate}

\subsection*{Data Structure Complexity}
\begin{itemize}[leftmargin=1.2em, itemsep=1pt]
    \item \textbf{List:} Index $\mathcal{O}(1)$, Append $\mathcal{O}(1)^*$, Insert/Del $\mathcal{O}(n)$.
    \item \textbf{Dict:} Key lookup, insert, delete average $\mathcal{O}(1)$.
    \item \textbf{Set:} Add, remove, lookup average $\mathcal{O}(1)$.
    \item \textbf{Deque:} Append/pop left and right $\mathcal{O}(1)$.
\end{itemize}

\subsection*{Dunder Methods Protocol}
\begin{itemize}[leftmargin=1.2em, itemsep=1pt]
    \item \texttt{\_\_init\_\_(self, ...)}: Initializer.
    \item \texttt{\_\_repr\_\_(self)}: Debug string.
    \item \texttt{\_\_len\_\_(self)}: Length (\texttt{int}).
    \item \texttt{\_\_getitem\_\_(self, k)}: Subscript indexing.
    \item \texttt{\_\_iter\_\_}, \texttt{\_\_next\_\_}: Iteration.
    \item \texttt{\_\_enter\_\_}, \texttt{\_\_exit\_\_}: Context manager.
\end{itemize}

\subsection*{NumPy Broadcasting}
Dimensions compatible if $d_A = d_B$ or one is $1$. Alignment evaluated trailing (right-to-left).
\begin{equation}
A(M, 1) \oplus B(1, N) \implies C(M, N)
\end{equation}

\subsection*{Pandas Indexing}
\begin{itemize}[leftmargin=1.2em, itemsep=1pt]
    \item \texttt{df.loc[lbl]}: Label-based (inclusive end).
    \item \texttt{df.iloc[pos]}: Integer-based (exclusive end).
\end{itemize}
\end{multicols}
\newpage


\tableofcontents
\newpage

\mainmatter

% SECTION 1: CONTINUOUS ASSESSMENT 1 (CA-1) --- UNITS 1 & 2
\chapter{Continuous Assessment 1 (CA-1) Suite}
\label{chap:ca1}
\noindent
Continuous Assessment 1 evaluates students on \textbf{Unit 1 (Python Environment Setup, Basics \& Computational Foundations)} and \textbf{Unit 2 (Control Flow, Functions, Problem-Solving \& Regular Expressions)}. Each paper contains \textbf{6 questions} carrying \textbf{6 marks each}, and students attempt any \textbf{5 questions} (Total = \textbf{30 Marks}, Time Allowed = \textbf{60 Minutes}).

\newpage
\section{Continuous Assessment 1 (CA-1) --- Set A}
\noindent\textbf{Course Code:} CSR101 \hfill \textbf{Max Marks:} 30\\
\textbf{Course Title:} Python Programming (AI \& Data Engg.) \hfill \textbf{Time Allowed:} 60 Minutes\\
\textbf{Instructions:} Answer any \textbf{5} out of the following 6 questions. Each carries \textbf{6 marks}.
\vspace{4mm}
\hrule
\vspace{4mm}

\noindent\textbf{Question 1 [Unit 1]:} Explain the internal structure of CPython's \texttt{PyObject}. What is the reference counting mechanism, and how does Python's cyclic garbage collector handle circular references?
\begin{solutionbox}
In CPython, every value allocated on the heap contains a \texttt{PyObject} header consisting of \texttt{ob\_refcnt} (64-bit reference counter) and \texttt{ob\_type} (pointer to type descriptor).
\begin{lstlisting}
import sys
x = [1, 2, 3]
print(sys.getrefcount(x))  # Prints 2 (x and getrefcount parameter)
\end{lstlisting}
\textbf{Cyclic Garbage Collection:}\\
Reference counting immediately reclaims objects when \texttt{ob\_refcnt == 0}. However, circular references ($A \to B \to A$) prevent the counter from reaching zero. CPython employs a generational cyclic garbage collector (\texttt{gc} module) that periodically tracks container objects (\texttt{list}, \texttt{dict}, user classes) across three generations (0, 1, 2). It isolates reference cycles by calculating unreachable subgraphs and destroys unreferenced cycles.
\end{solutionbox}
\begin{markingrubric}
\textbf{Marking Scheme:} \texttt{PyObject} struct fields (2M); Reference counting increment/decrement lifecycle (2M); Generational cyclic collection of isolated loops (2M).
\end{markingrubric}

\noindent\textbf{Question 2 [Unit 1]:} Discuss the difference between integer division \texttt{//} and modulo \texttt{\%} in Python when negative operands are involved. Evaluate \texttt{-25 // 4} and \texttt{-25 \% 4}.
\begin{solutionbox}
Python enforces the mathematical floor division invariant $a = (a // b) \cdot b + (a \% b)$.
For $a = -25, b = 4$:
\begin{equation}
-25 / 4 = -6.25 \implies \lfloor -6.25 \rfloor = -7
\end{equation}
Then the remainder is:
\begin{equation}
r = -25 - (-7 \cdot 4) = -25 - (-28) = 3
\end{equation}
Thus \texttt{-25 // 4 == -7} and \texttt{-25 \% 4 == 3}.
\end{solutionbox}
\begin{markingrubric}
\textbf{Marking Scheme:} Floor definition towards $-\infty$ (2M); Calculation of quotient $-7$ (2M); Derivation of modulo $+3$ (2M).
\end{markingrubric}

\noindent\textbf{Question 3 [Unit 2]:} Implement a function \texttt{extract\_emails(text)} using regular expressions that extracts and returns all unique valid email addresses sorted alphabetically.
\begin{solutionbox}
\begin{lstlisting}
import re

def extract_emails(text: str) -> list:
    pattern = re.compile(r'\b[A-Za-z0-9._%+-]+@[A-Za-z0-9.-]+\.[A-Z|a-z]{2,}\b')
    emails = pattern.findall(text)
    # Return unique sorted emails
    return sorted(list(set(emails)))
\end{lstlisting}
\end{solutionbox}
\begin{markingrubric}
\textbf{Marking Scheme:} Correct regex pattern syntax (3M); Set deduplication and sorting logic (2M); Function signature and PEP 8 compliance (1M).
\end{markingrubric}

\noindent\textbf{Question 4 [Unit 2]:} Write a recursive Python function \texttt{binary\_search(arr, target, low, high)} that returns the index of target in a sorted list or $-1$ if absent. State the recurrence relation and space complexity.
\begin{solutionbox}
\begin{lstlisting}
def binary_search(arr, target, low, high):
    if low > high:
        return -1
    mid = (low + high) // 2
    if arr[mid] == target:
        return mid
    elif arr[mid] > target:
        return binary_search(arr, target, low, mid - 1)
    else:
        return binary_search(arr, target, mid + 1, high)
\end{lstlisting}
\textbf{Recurrence Relation:} $T(n) = T(n/2) + \mathcal{O}(1) \implies \mathcal{O}(\log n)$ time. Space complexity: $\mathcal{O}(\log n)$ stack frames.
\end{solutionbox}
\begin{markingrubric}
\textbf{Marking Scheme:} Base cases and mid index calculation (2M); Recursive calls logic (2M); Recurrence relation \& space complexity (2M).
\end{markingrubric}

\noindent\textbf{Question 5 [Unit 2]:} Explain variable scoping and the LEGB rule. Demonstrate the difference between \texttt{global} and \texttt{nonlocal} keywords using a clear code example.
\begin{solutionbox}
\begin{lstlisting}
x = "Global"

def outer():
    x = "Enclosing"
    def inner_nonlocal():
        nonlocal x
        x = "Modified Enclosing"
    def inner_global():
        global x
        x = "Modified Global"
    inner_nonlocal()
    print("Outer after nonlocal:", x)  # Prints: Modified Enclosing
outer()
print("Global after outer:", x)        # Prints: Global
\end{lstlisting}
\end{solutionbox}
\begin{markingrubric}
\textbf{Marking Scheme:} Explanation of LEGB 4 levels (2M); Correct \texttt{nonlocal} usage and behavior (2M); Correct \texttt{global} demonstration (2M).
\end{markingrubric}

\noindent\textbf{Question 6 [Unit 2]:} Explain why using a mutable object (like a list or dict) as a default argument in a Python function is an anti-pattern. Provide the standard design solution.
\begin{solutionbox}
Default arguments are evaluated once when the function definition is executed, creating a single shared object across all invocations.
\begin{lstlisting}
# Standard Idiom: None Sentinel
def append_item(val, container=None):
    if container is None:
        container = []
    container.append(val)
    return container
\end{lstlisting}
\end{solutionbox}
\begin{markingrubric}
\textbf{Marking Scheme:} Identification of definition-time evaluation (3M); Correct \texttt{None} sentinel pattern implementation (3M).
\end{markingrubric}

\newpage
\section{Continuous Assessment 1 (CA-1) --- Set B}
\noindent\textbf{Course Code:} CSR101 \hfill \textbf{Max Marks:} 30\\
\textbf{Course Title:} Python Programming (AI \& Data Engg.) \hfill \textbf{Time Allowed:} 60 Minutes\\
\textbf{Instructions:} Answer any \textbf{5} out of the following 6 questions. Each carries \textbf{6 marks}.
\vspace{4mm}
\hrule
\vspace{4mm}

\noindent\textbf{Question 1 [Unit 1]:} What is PEP 8? Discuss four crucial guidelines prescribed by PEP 8 concerning indentation, variable naming conventions, imports, and maximum line length.
\begin{solutionbox}
PEP 8 is the official Python Enhancement Proposal establishing the style guide for Python code.
\begin{enumerate}[leftmargin=1.5em, itemsep=2pt]
    \item \textbf{Indentation:} 4 spaces per indentation level. Never mix tabs and spaces.
    \item \textbf{Naming Conventions:} \texttt{snake\_case} for functions and variables, \texttt{PascalCase} for classes, \texttt{ALL\_CAPS} for constants.
    \item \textbf{Imports:} Separate lines at top of file, grouped in 3 sections: Standard library, third-party libraries, local application imports.
    \item \textbf{Line Length:} Maximum limit of 79 characters for code and 72 characters for docstrings/comments.
\end{enumerate}
\end{solutionbox}
\begin{markingrubric}
\textbf{Marking Scheme:} Definition of PEP 8 (1M); Explanation of 4 rules with correct syntax examples (5M).
\end{markingrubric}

\noindent\textbf{Question 2 [Unit 1]:} Explain the difference between identity operator \texttt{is} and equality operator \texttt{==}. When can \texttt{a is b} evaluate to \texttt{True} for two integer variables?
\begin{solutionbox}
$\bullet$ \texttt{==} compares \textbf{values} by calling \texttt{a.\_\_eq\_\_(b)}.\\
$\bullet$ \texttt{is} compares \textbf{memory identities} (\texttt{id(a) == id(b)}).\\
For integers, \texttt{a is b} is True when both refer to an interned integer in the range $[-5, 256]$.
\end{solutionbox}
\begin{markingrubric}
\textbf{Marking Scheme:} Difference between value equality and memory identity (3M); Explanation of CPython integer interning cache $[-5, 256]$ (3M).
\end{markingrubric}

\noindent\textbf{Question 3 [Unit 2]:} Implement a function \texttt{flatten\_nested\_list(nested)} using recursion that takes an arbitrarily deeply nested list of integers and returns a flat 1D list.
\begin{solutionbox}
\begin{lstlisting}
def flatten_nested_list(nested):
    flat = []
    for item in nested:
        if isinstance(item, list):
            flat.extend(flatten_nested_list(item))
        else:
            flat.append(item)
    return flat
\end{lstlisting}
\end{solutionbox}
\begin{markingrubric}
\textbf{Marking Scheme:} Base case detection via \texttt{isinstance} (2M); Recursive flattening and list concatenation (3M); Testing with empty/deep inputs (1M).
\end{markingrubric}

\noindent\textbf{Question 4 [Unit 2]:} Write a Python function that uses the \texttt{re} module to validate a strong password (minimum 8 characters, at least 1 uppercase, 1 lowercase, 1 digit, and 1 special character).
\begin{solutionbox}
\begin{lstlisting}
import re

def validate_password(pwd: str) -> bool:
    pattern = re.compile(r'^(?=.*[a-z])(?=.*[A-Z])(?=.*\d)(?=.*[@$!%*?&])[A-Za-z\d@$!%*?&]{8,}$')
    return bool(pattern.match(pwd))
\end{lstlisting}
\end{solutionbox}
\begin{markingrubric}
\textbf{Marking Scheme:} Use of lookahead assertions \texttt{(?=...)} (4M); Character set and length boundary constraints (2M).
\end{markingrubric}

\noindent\textbf{Question 5 [Unit 2]:} Explain the loop \texttt{else} block invariant in Python. Give a worked example demonstrating how loop \texttt{else} avoids boolean flag variables.
\begin{solutionbox}
The \texttt{else} clause of a loop executes if and only if the loop completes normally without encountering a \texttt{break} statement.
\begin{lstlisting}
def find_item(target, items):
    for item in items:
        if item == target:
            print("Item found!")
            break
    else:
        print("Target not found in collection.")
\end{lstlisting}
\end{solutionbox}
\begin{markingrubric}
\textbf{Marking Scheme:} Exact condition triggering \texttt{else} block (3M); Code implementation without boolean flag (3M).
\end{markingrubric}

\noindent\textbf{Question 6 [Unit 2]:} Describe Python's \texttt{*args} and \texttt{**kwargs} mechanics. Write a function \texttt{pipeline\_logger(*args, **kwargs)} that prints formatted positional and keyword parameters.
\begin{solutionbox}
\begin{lstlisting}
def pipeline_logger(*args, **kwargs):
    print(f"Positional arguments ({len(args)}): {args}")
    for k, v in kwargs.items():
        print(f"Keyword argument: {k} -> {v}")
\end{lstlisting}
\end{solutionbox}
\begin{markingrubric}
\textbf{Marking Scheme:} Tuple packing for \texttt{*args} (2M); Dict packing for \texttt{**kwargs} (2M); Code implementation (2M).
\end{markingrubric}


% SECTION 2: MID-TERM EXAMINATION SUITE --- UNITS 1, 2 & 3
\chapter{Mid-Term Examination Suite (Units 1, 2 \& 3)}
\label{chap:midterm}
\noindent
The Mid-Term Examination evaluates students across \textbf{Unit 1 (Basics)}, \textbf{Unit 2 (Functions \& Logic)}, and \textbf{Unit 3 (Data Structures \& OOP)}. Each paper consists of \textbf{15 Multiple-Choice Questions (MCQs)} carrying \textbf{1 mark each} with a negative marking deduction of \textbf{0.25 marks} per incorrect answer, plus \textbf{3 Subjective Coding \& Design Problems} carrying \textbf{5 marks each} (Total = \textbf{30 Marks}, Time Allowed = \textbf{90 Minutes}).

\newpage
\section{Mid-Term Examination --- Set A (Units 1, 2 \& 3)}
\noindent\textbf{Course Code:} CSR101 \hfill \textbf{Max Marks:} 30\\
\textbf{Course Title:} Python Programming (AI \& Data Engg.) \hfill \textbf{Time Allowed:} 90 Minutes\\
\textbf{Rules:} Part A contains 15 MCQs (1 Mark each, $-0.25$ negative marking). Part B contains 3 Subjective Coding problems of 5 Marks each.
\vspace{4mm}
\hrule
\vspace{4mm}

\subsection*{Part A: Multiple Choice Questions (15 Marks)}
\begin{enumerate}[label=\textbf{\arabic*.}]
    \item Which tool is commonly used to write and run Python code? \hfill \textbf{[Ans: Jupyter Notebook / VS Code]}
    \item What is the time complexity of appending an element to a Python list? \hfill \textbf{[Ans: Amortized $\mathcal{O}(1)$]}
    \item What is printed by \texttt{print(type(lambda x: x))}? \hfill \textbf{[Ans: \texttt{<class 'function'>}]}
    \item Which data structure uses hash tables in CPython? \hfill \textbf{[Ans: Dictionary and Set]}
    \item What does \texttt{re.match(r'\textbackslash d+', 'abc123')} return? \hfill \textbf{[Ans: \texttt{None} (matches only at start)]}
    \item Which method is invoked by the built-in \texttt{len()} function? \hfill \textbf{[Ans: \texttt{\_\_len\_\_()}]}
    \item What is the value of \texttt{bool([])}? \hfill \textbf{[Ans: \texttt{False}]}
    \item What does the \texttt{nonlocal} keyword do? \hfill \textbf{[Ans: Binds identifier to nearest enclosing scope]}
    \item Which algorithm resolves Python's Method Resolution Order (MRO)? \hfill \textbf{[Ans: C3 Linearization]}
    \item What is the result of \texttt{2 ** 3 ** 2}? \hfill \textbf{[Ans: 512 (Right-to-left associativity)]}
    \item How does Python handle name mangling for \texttt{\_\_x} in class \texttt{Foo}? \hfill \textbf{[Ans: \texttt{\_Foo\_\_x}]}
    \item What is returned by a generator function containing \texttt{yield}? \hfill \textbf{[Ans: A generator iterator object]}
    \item What is the default return value of a function lacking a \texttt{return} statement? \hfill \textbf{[Ans: \texttt{None}]}
    \item In CPython, which integers are cached as singletons? \hfill \textbf{[Ans: $[-5, 256]$]}
    \item Which operator performs set intersection? \hfill \textbf{[Ans: \texttt{\&}]}
\end{enumerate}

\subsection*{Part B: Subjective Coding \& Design Problems (15 Marks)}

\noindent\textbf{Question 16 [Unit 3, 5 Marks]:} Design a Python class \texttt{Polynomial} representing univariate polynomials using a dictionary of $\{ \text{degree}: \text{coefficient} \}$. Implement \texttt{\_\_add\_\_} and \texttt{\_\_call\_\_} (evaluation at $x$).
\begin{solutionbox}
\begin{lstlisting}
class Polynomial:
    def __init__(self, terms: dict):
        self.terms = {deg: coef for deg, coef in terms.items() if coef != 0}

    def __add__(self, other):
        res = self.terms.copy()
        for deg, coef in other.terms.items():
            res[deg] = res.get(deg, 0) + coef
        return Polynomial(res)

    def __call__(self, x: float) -> float:
        return sum(coef * (x ** deg) for deg, coef in self.terms.items())
\end{lstlisting}
\end{solutionbox}
\begin{markingrubric}
\textbf{Marking Scheme:} Data representation in dictionary (1M); Implementation of \texttt{\_\_add\_\_} operator overload (2M); Implementation of \texttt{\_\_call\_\_} evaluation (2M).
\end{markingrubric}

\noindent\textbf{Question 17 [Unit 2, 5 Marks]:} Implement a memoization decorator \texttt{@memoize} from scratch without using \texttt{functools.lru\_cache} and demonstrate it on the Fibonacci sequence.
\begin{solutionbox}
\begin{lstlisting}
def memoize(func):
    cache = {}
    def wrapper(*args):
        if args not in cache:
            cache[args] = func(*args)
        return cache[args]
    return wrapper

@memoize
def fib(n):
    if n < 2: return n
    return fib(n - 1) + fib(n - 2)
\end{lstlisting}
\textbf{Complexity Improvement:} Transforms naive Fibonacci from $\mathcal{O}(2^n)$ to $\mathcal{O}(n)$ time.
\end{solutionbox}
\begin{markingrubric}
\textbf{Marking Scheme:} Closure with cache dictionary (2M); Key check and conditional computation (2M); Asymptotic complexity explanation (1M).
\end{markingrubric}

\noindent\textbf{Question 18 [Unit 3, 5 Marks]:} Explain C3 Linearization. Compute the MRO for the inheritance hierarchy: class A, class B(A), class C(A), class D(B, C).
\begin{solutionbox}
The MRO formula is $L(D) = [D] + \operatorname{merge}(L(B), L(C), [B, C])$.
\begin{align*}
L(A) &= [A, O] \\
L(B) &= [B, A, O] \\
L(C) &= [C, A, O] \\
L(D) &= [D] + \operatorname{merge}([B, A, O], [C, A, O], [B, C]) \\
     &= [D, B] + \operatorname{merge}([A, O], [C, A, O], [C]) \\
     &= [D, B, C] + \operatorname{merge}([A, O], [A, O]) \\
     &= [D, B, C, A, O]
\end{align*}
Result: \texttt{D -> B -> C -> A -> object}.
\end{solutionbox}
\begin{markingrubric}
\textbf{Marking Scheme:} Statement of C3 merge algorithm (2M); Step-by-step resolution of candidate lists (2M); Final correct MRO sequence (1M).
\end{markingrubric}

\newpage
\section{Mid-Term Examination --- Set B (Units 1, 2 \& 3)}
\noindent\textbf{Course Code:} CSR101 \hfill \textbf{Max Marks:} 30\\
\textbf{Course Title:} Python Programming (AI \& Data Engg.) \hfill \textbf{Time Allowed:} 90 Minutes\\
\vspace{4mm}
\hrule
\vspace{4mm}

\subsection*{Part A: Multiple Choice Questions (15 Marks)}
\begin{enumerate}[label=\textbf{\arabic*.}]
    \item What is the output of \texttt{bool("False")}? \hfill \textbf{[Ans: \texttt{True} (non-empty string)]}
    \item Which built-in function returns an iterator yielding $(index, value)$ pairs? \hfill \textbf{[Ans: \texttt{enumerate()}]}
    \item How are keyword-only arguments denoted in a parameter list? \hfill \textbf{[Ans: After a bare \texttt{*} or \texttt{*args}]}
    \item What is the space complexity of a generator expression compared to a list comprehension? \hfill \textbf{[Ans: $\mathcal{O}(1)$ vs $\mathcal{O}(n)$]}
    \item In a class, how are class methods declared? \hfill \textbf{[Ans: With the \texttt{@classmethod} decorator and \texttt{cls} parameter]}
    \item What is the output of \texttt{[1, 2] * 2}? \hfill \textbf{[Ans: \texttt{[1, 2, 1, 2]}]}
    \item What is printed by \texttt{set([1, 2, 2, 3])}? \hfill \textbf{[Ans: \texttt{\{1, 2, 3\}}]}
    \item What is the purpose of \texttt{\_\_slots\_\_}? \hfill \textbf{[Ans: Prevents dynamic \texttt{\_\_dict\_\_} creation, saving memory]}
    \item What does \texttt{re.split(r'\textbackslash s+', text)} do? \hfill \textbf{[Ans: Splits string on one or more whitespace characters]}
    \item Can a tuple containing a list be used as a dictionary key? \hfill \textbf{[Ans: No, raises \texttt{TypeError: unhashable type: 'list'}]}
    \item What is the output of \texttt{10 // -3}? \hfill \textbf{[Ans: \texttt{-4}]}
    \item What does \texttt{zip([1, 2], ['a', 'b'])} return? \hfill \textbf{[Ans: Iterator yielding tuples: \texttt{(1, 'a'), (2, 'b')}]}
    \item Which dunder method is called when an object is used in a \texttt{with} statement? \hfill \textbf{[Ans: \texttt{\_\_enter\_\_()}]}
    \item What does \texttt{isinstance(True, int)} evaluate to? \hfill \textbf{[Ans: \texttt{True} (\texttt{bool} subclasses \texttt{int})]}
    \item What does \texttt{re.sub(r'cat', 'dog', text)} do? \hfill \textbf{[Ans: Replaces all occurrences of 'cat' with 'dog']}
\end{enumerate}

\subsection*{Part B: Subjective Coding \& Design Problems (15 Marks)}

\noindent\textbf{Question 16 [Unit 3, 5 Marks]:} Implement a custom iterator class \texttt{PrimeIterator(limit)} that yields consecutive prime numbers up to \texttt{limit} implementing the \texttt{\_\_iter\_\_} and \texttt{\_\_next\_\_} protocol.
\begin{solutionbox}
\begin{lstlisting}
class PrimeIterator:
    def __init__(self, limit: int):
        self.limit = limit
        self.current = 2

    def __iter__(self):
        return self

    def _is_prime(self, n: int) -> bool:
        for d in range(2, int(n**0.5) + 1):
            if n % d == 0: return False
        return True

    def __next__(self):
        while self.current <= self.limit:
            n = self.current
            self.current += 1
            if self._is_prime(n):
                return n
        raise StopIteration
\end{lstlisting}
\end{solutionbox}
\begin{markingrubric}
\textbf{Marking Scheme:} Class initialization and \texttt{\_\_iter\_\_} returning self (1M); Prime checking logic (2M); \texttt{\_\_next\_\_} state mutation and \texttt{StopIteration} raising (2M).
\end{markingrubric}

\noindent\textbf{Question 17 [Unit 3, 5 Marks]:} Implement a custom Vector class in 3D space supporting vector addition (\texttt{+}), dot product (\texttt{@}), and magnitude representation.
\begin{solutionbox}
\begin{lstlisting}
import math

class Vector3D:
    def __init__(self, x: float, y: float, z: float):
        self.x, self.y, self.z = x, y, z

    def __add__(self, other):
        return Vector3D(self.x + other.x, self.y + other.y, self.z + other.z)

    def __matmul__(self, other):
        return self.x * other.x + self.y * other.y + self.z * other.z

    def magnitude(self) -> float:
        return math.sqrt(self.x**2 + self.y**2 + self.z**2)

    def __repr__(self):
        return f"Vector3D({self.x}, {self.y}, {self.z})"
\end{lstlisting}
\end{solutionbox}
\begin{markingrubric}
\textbf{Marking Scheme:} \texttt{\_\_add\_\_} operator overload (1.5M); \texttt{\_\_matmul\_\_} dot product (1.5M); Magnitude and representation (2M).
\end{markingrubric}

\noindent\textbf{Question 18 [Unit 2, 5 Marks]:} Write a generator function \texttt{sliding\_window(iterable, k)} that yields consecutive sliding windows of length $k$ from a sequence.
\begin{solutionbox}
\begin{lstlisting}
from collections import deque

def sliding_window(iterable, k: int):
    it = iter(iterable)
    window = deque(maxlen=k)
    for _ in range(k):
        try:
            window.append(next(it))
        except StopIteration:
            return
    yield tuple(window)
    for elem in it:
        window.append(elem)
        yield tuple(window)
\end{lstlisting}
\end{solutionbox}
\begin{markingrubric}
\textbf{Marking Scheme:} Proper use of generator \texttt{yield} (2M); Efficient $\mathcal{O}(1)$ window updates using \texttt{deque} (2M); Edge cases handling (1M).
\end{markingrubric}


% SECTION 3: CONTINUOUS ASSESSMENT 2 (CA-2) --- UNITS 4 & 5
\chapter{Continuous Assessment 2 (CA-2) Suite}
\label{chap:ca2}
\noindent
Continuous Assessment 2 evaluates students on \textbf{Unit 4 (Modules, Packages \& Exception Handling)} and \textbf{Unit 5 (Matrices, Vectorized Computing \& NumPy Linear Algebra)}. Each paper contains \textbf{6 questions} carrying \textbf{6 marks each}, and students attempt any \textbf{5 questions} (Total = \textbf{30 Marks}, Time Allowed = \textbf{60 Minutes}).

\newpage
\section{Continuous Assessment 2 (CA-2) --- Set A}
\noindent\textbf{Course Code:} CSR101 \hfill \textbf{Max Marks:} 30\\
\textbf{Course Title:} Python Programming (AI \& Data Engg.) \hfill \textbf{Time Allowed:} 60 Minutes\\
\textbf{Instructions:} Answer any \textbf{5} out of the following 6 questions (Units 4 \& 5). Each carries \textbf{6 marks}.
\vspace{4mm}
\hrule
\vspace{4mm}

\noindent\textbf{Question 1 [Unit 4]:} Explain Python's module import resolution algorithm. What steps does CPython follow from the moment \texttt{import mymodule} is executed?
\begin{solutionbox}
1. \textbf{Module Cache Check:} Checks \texttt{sys.modules['mymodule']}. If cached, returns reference.\\
2. \textbf{Built-in Modules:} Searches \texttt{sys.builtin\_module\_names}.\\
3. \textbf{Search in \texttt{sys.path}:} Searches script directory, \texttt{PYTHONPATH}, standard library, and site-packages.\\
4. \textbf{Compilation:} Compiles source to bytecode (\texttt{.pyc}) if not already updated.\\
5. \textbf{Execution:} Executes top-level module code in a fresh module namespace and stores reference in \texttt{sys.modules}.
\end{solutionbox}
\begin{markingrubric}
\textbf{Marking Scheme:} Cache check in \texttt{sys.modules} (1.5M); Built-in lookup (1.5M); \texttt{sys.path} directories search (1.5M); Module execution and caching (1.5M).
\end{markingrubric}

\noindent\textbf{Question 2 [Unit 4]:} Design a custom exception class \texttt{InvalidModelConfigurationError} with attributes \texttt{model\_name} and \texttt{param\_name}. Show how to raise and catch it with custom error messages.
\begin{solutionbox}
\begin{lstlisting}
class InvalidModelConfigurationError(Exception):
    def __init__(self, model_name: str, param_name: str, message: str = "Invalid config parameter"):
        self.model_name = model_name
        self.param_name = param_name
        super().__init__(f"[{model_name}] {message}: '{param_name}'")

# Demonstration
try:
    raise InvalidModelConfigurationError("Transformer-XL", "learning_rate", "Must be positive")
except InvalidModelConfigurationError as err:
    print(f"Handled error: {err}")
\end{lstlisting}
\end{solutionbox}
\begin{markingrubric}
\textbf{Marking Scheme:} Subclassing \texttt{Exception} and \texttt{\_\_init\_\_} super call (3M); Custom attributes and formatted message (2M); Try-except handling syntax (1M).
\end{markingrubric}

\noindent\textbf{Question 3 [Unit 5]:} What is NumPy broadcasting? State the exact conditions required for two arrays of arbitrary dimensions to be broadcast-compatible. Illustrate with an example.
\begin{solutionbox}
Broadcasting allows element-wise operations between arrays of different shapes without copying data.
\textbf{Conditions (evaluated right-to-left):}\\
1. Dimensions are equal, OR\\
2. One of the dimensions is 1.\\
\textbf{Example:} Array $A$ of shape $(4, 1, 8)$ and Array $B$ of shape $(5, 8)$ broadcast to result shape $(4, 5, 8)$.
\end{solutionbox}
\begin{markingrubric}
\textbf{Marking Scheme:} Broadcasting definition and zero-copy concept (2M); Two compatibility conditions (2M); Concrete shape alignment example (2M).
\end{markingrubric}

\noindent\textbf{Question 4 [Unit 5]:} Implement a vectorized Python function \texttt{pairwise\_euclidean\_distance(X, Y)} using pure NumPy broadcasting without Python loops.
\begin{solutionbox}
\begin{lstlisting}
import numpy as np

def pairwise_euclidean_distance(X: np.ndarray, Y: np.ndarray) -> np.ndarray:
    # X shape: (M, D) -> reshape to (M, 1, D)
    # Y shape: (N, D) -> reshape to (1, N, D)
    diff = X[:, np.newaxis, :] - Y[np.newaxis, :, :]
    # diff shape: (M, N, D)
    return np.sqrt(np.sum(diff ** 2, axis=-1))
\end{lstlisting}
\end{solutionbox}
\begin{markingrubric}
\textbf{Marking Scheme:} Dimension expansion via \texttt{np.newaxis} (3M); Vectorized difference and squared summation (2M); Return shape $(M, N)$ (1M).
\end{markingrubric}

\noindent\textbf{Question 5 [Unit 5]:} Explain Cholesky decomposition ($A = L L^T$). Why is it preferred over standard Gaussian elimination for symmetric positive-definite linear systems?
\begin{solutionbox}
Cholesky decomposes a symmetric positive-definite matrix $A$ into a lower triangular matrix $L$ such that $A = L L^T$.
\textbf{Advantages:}\\
1. \textbf{Efficiency:} Requires $\frac{1}{3} n^3$ floating-point operations compared to $\frac{2}{3} n^3$ for general LU factorization (twice as fast).\\
2. \textbf{Numerical Stability:} No pivoting is required, eliminating row swapping overhead.
\end{solutionbox}
\begin{markingrubric}
\textbf{Marking Scheme:} Mathematical definition and triangular form (2M); Complexity comparison ($n^3/3$ vs $2n^3/3$) (2M); Numerical stability without pivoting (2M).
\end{markingrubric}

\noindent\textbf{Question 6 [Unit 4]:} What is a context manager? Implement a custom context manager class \texttt{ExecutionTimer} that measures and logs the elapsed time of any code block.
\begin{solutionbox}
\begin{lstlisting}
import time

class ExecutionTimer:
    def __enter__(self):
        self.start = time.perf_counter()
        return self

    def __exit__(self, exc_type, exc_val, exc_tb):
        elapsed = time.perf_counter() - self.start
        print(f"Code block executed in {elapsed:.6f} seconds.")
        return False  # Do not suppress exceptions
\end{lstlisting}
\end{solutionbox}
\begin{markingrubric}
\textbf{Marking Scheme:} Definition of PEP 343 context protocol (2M); Implementation of \texttt{\_\_enter\_\_} (2M); Implementation of \texttt{\_\_exit\_\_} and exception propagation (2M).
\end{markingrubric}

\newpage
\section{Continuous Assessment 2 (CA-2) --- Set B}
\noindent\textbf{Course Code:} CSR101 \hfill \textbf{Max Marks:} 30\\
\textbf{Course Title:} Python Programming (AI \& Data Engg.) \hfill \textbf{Time Allowed:} 60 Minutes\\
\vspace{4mm}
\hrule
\vspace{4mm}

\noindent\textbf{Question 1 [Unit 4]:} Differentiate between compile-time syntax errors and runtime exceptions. Describe the role of the \texttt{finally} block in enterprise Python applications.
\begin{solutionbox}
$\bullet$ \textbf{Syntax Errors:} Detected during AST parsing before bytecode execution starts.\\
$\bullet$ \textbf{Runtime Exceptions:} Raised during execution when invalid operations occur (e.g., \texttt{ZeroDivisionError}, \texttt{IndexError}).\\
$\bullet$ \textbf{\texttt{finally} block:} Guarantees cleanup execution (closing database connections, releasing mutexes) regardless of whether an exception occurred, was caught, or propagated.
\end{solutionbox}
\begin{markingrubric}
\textbf{Marking Scheme:} Syntax vs Runtime distinction (3M); Guarantee semantics of \texttt{finally} block (3M).
\end{markingrubric}

\noindent\textbf{Question 2 [Unit 4]:} How does Python distinguish between importing a module and running it directly? Explain the role of \texttt{\_\_name\_\_ == "\_\_main\_\_"}.
\begin{solutionbox}
CPython sets the built-in variable \texttt{\_\_name\_\_} to \texttt{"\_\_main\_\_"} when the script is executed directly from the command line. When the file is imported by another module, \texttt{\_\_name\_\_} is set to the module's filename. This guard allows a file to serve as both an importable library and an executable CLI tool.
\end{solutionbox}
\begin{markingrubric}
\textbf{Marking Scheme:} Explanation of \texttt{\_\_name\_\_} assignment (3M); Standalone vs import usage pattern (3M).
\end{markingrubric}

\noindent\textbf{Question 3 [Unit 5]:} Explain array strides in NumPy. How does NumPy perform matrix transposition without allocating new memory?
\begin{solutionbox}
Strides represent the byte step necessary to advance one index along an axis. For an array of shape $(M, N)$ with 8-byte elements, default C-strides are $(8N, 8)$. Transposition simply swaps the strides to $(8, 8N)$ and shape to $(N, M)$. This metadata-only operation runs in $\mathcal{O}(1)$ time with zero data copying.
\end{solutionbox}
\begin{markingrubric}
\textbf{Marking Scheme:} Stride definition and byte math (3M); Transposition as metadata stride swap (3M).
\end{markingrubric}

\noindent\textbf{Question 4 [Unit 5]:} Solve the linear system $A x = b$ using NumPy, where $A = \begin{bmatrix} 4 & 2 \\ 2 & 5 \end{bmatrix}$ and $b = \begin{bmatrix} 8 \\ 10 \end{bmatrix}$. Compute the condition number of $A$.
\begin{solutionbox}
\begin{lstlisting}
import numpy as np

A = np.array([[4.0, 2.0], [2.0, 5.0]])
b = np.array([8.0, 10.0])

x = np.linalg.solve(A, b)
cond = np.linalg.cond(A)
print(f"Solution x: {x}, Condition Number: {cond:.4f}")
\end{lstlisting}
Analytical solution: $\det(A) = 20 - 4 = 16$. $x_1 = (40 - 20)/16 = 1.25$, $x_2 = (40 - 16)/16 = 1.5$.
\end{solutionbox}
\begin{markingrubric}
\textbf{Marking Scheme:} Matrix setup and solve API (3M); Condition number interpretation (3M).
\end{markingrubric}

\noindent\textbf{Question 5 [Unit 5]:} Compare standard dense arrays with SciPy sparse matrices (\texttt{csr\_matrix}). In what AI scenarios are sparse matrices essential?
\begin{solutionbox}
Dense arrays store every value, requiring $\mathcal{O}(M \times N)$ memory. Compressed Sparse Row (CSR) stores only non-zero values using three arrays: \texttt{data}, \texttt{indices}, and \texttt{indptr}, saving gigabytes of memory.
\textbf{Applications:} NLP TF-IDF matrices, graph adjacency matrices, recommendation collaborative filtering.
\end{solutionbox}
\begin{markingrubric}
\textbf{Marking Scheme:} Memory representation comparison (3M); CSR structure and AI applications (3M).
\end{markingrubric}

\noindent\textbf{Question 6 [Unit 4]:} What is exception chaining? Explain the difference between \texttt{raise new\_exc from original\_exc} and \texttt{raise new\_exc from None}.
\begin{solutionbox}
$\bullet$ \texttt{raise new\_exc from original\_exc}: Explicitly links exceptions, setting \texttt{\_\_cause\_\_} and preserving root-cause traceback.\\
$\bullet$ \texttt{raise new\_exc from None}: Suppresses the context, setting \texttt{\_\_suppress\_context\_\_ = True} to hide internal implementation details from end users.
\end{solutionbox}
\begin{markingrubric}
\textbf{Marking Scheme:} Explanation of \texttt{\_\_cause\_\_} chaining (3M); Use cases for suppressing context with \texttt{from None} (3M).
\end{markingrubric}


% SECTION 4: END-TERM EXAMINATION SUITE --- UNITS 1 TO 6
\chapter{End-Term Examination Suite (Units 1 to 6 Comprehensive --- 70 Marks)}
\label{chap:endterm}
\noindent
The End-Term Examination is the comprehensive capstone assessment evaluating all six units of the CSR101 curriculum. Each paper consists of \textbf{Part A} (15 Objective Questions of 1 Mark each), \textbf{Part B} (5 Analytical \& Coding Questions of 5 Marks each), and \textbf{Part C} (3 In-Depth Engineering \& Pipeline Design Problems of 10 Marks each) (Total = \textbf{70 Marks}, Time Allowed = \textbf{3 Hours}).

\newpage
\section{End-Term Comprehensive Examination --- Set A (Units 1 to 6)}
\noindent\textbf{Course Code:} CSR101 \hfill \textbf{Max Marks:} 70\\
\textbf{Course Title:} Python Programming (AI \& Data Engg.) \hfill \textbf{Time Allowed:} 3 Hours\\
\textbf{Structure:} Part A: 15 Objective Questions (15 Marks); Part B: 5 Analytical Questions (25 Marks); Part C: 3 Comprehensive Pipeline Design Problems (30 Marks).
\vspace{4mm}
\hrule
\vspace{4mm}

\subsection*{Part A: Objective Conceptual Questions (15 Marks)}
\begin{enumerate}[label=\textbf{\arabic*.}]
    \item What is the role of \texttt{\_\_pycache\_\_} in Python projects? \hfill \textbf{[Ans: Stores compiled bytecode files (\texttt{.pyc})]}
    \item What is the output of \texttt{type(1 / 1)}? \hfill \textbf{[Ans: \texttt{<class 'float'>}]}
    \item Which regex quantifier matches between 2 and 5 occurrences? \hfill \textbf{[Ans: \texttt{\{2,5\}}]}
    \item What is the difference between \texttt{list.append()} and \texttt{list.extend()}? \hfill \textbf{[Ans: \texttt{append} adds object as single item; \texttt{extend} iterates]}
    \item Can a Python function return multiple values? How? \hfill \textbf{[Ans: Yes, packed automatically into an immutable tuple]}
    \item What does \texttt{re.findall(r'\textbackslash d+', 'h3ll0 42')} return? \hfill \textbf{[Ans: \texttt{['3', '0', '42']}]}
    \item What is the time complexity of searching a key in a dict? \hfill \textbf{[Ans: Average $\mathcal{O}(1)$, worst-case $\mathcal{O}(n)$]}
    \item What is the difference between a shallow copy and a deep copy? \hfill \textbf{[Ans: Shallow copies outer container; deep recursively clones nested objects]}
    \item What is the purpose of \texttt{sys.path}? \hfill \textbf{[Ans: List of directory paths searched during module imports]}
    \item What exception is raised when dividing by zero in floating point math? \hfill \textbf{[Ans: \texttt{ZeroDivisionError}]}
    \item In NumPy, what attribute returns the tuple of array dimensions? \hfill \textbf{[Ans: \texttt{shape}]}
    \item What does \texttt{np.eye(3)} create? \hfill \textbf{[Ans: A $3 \times 3$ identity matrix]}
    \item In Pandas, how does \texttt{.loc} differ from \texttt{.iloc}? \hfill \textbf{[Ans: \texttt{.loc} is label-based; \texttt{.iloc} is integer position-based]}
    \item Which file mode opens a file for writing binary data, truncating existing content? \hfill \textbf{[Ans: \texttt{'wb'}]}
    \item What method saves an active Matplotlib figure to disk? \hfill \textbf{[Ans: \texttt{plt.savefig()}]}
\end{enumerate}

\subsection*{Part B: Analytical \& Coding Problems (25 Marks)}

\noindent\textbf{Question 16 [5 Marks]:} Explain CPython's memory management model, distinguishing between stack frames, heap allocation, and the GIL (Global Interpreter Lock).
\begin{solutionbox}
$\bullet$ \textbf{Stack Frames:} Store local variable references, instruction pointers, and evaluation stacks for active function invocations.\\
$\bullet$ \textbf{Heap Allocation:} All Python objects (\texttt{PyObject}) reside on the private heap managed by PyMalloc.\\
$\bullet$ \textbf{GIL (Global Interpreter Lock):} A mutex preventing multiple native OS threads from executing CPython bytecode simultaneously, ensuring thread-safe reference counting.
\end{solutionbox}
\begin{markingrubric}
\textbf{Marking Scheme:} Stack frame execution role (1.5M); Heap allocation and PyMalloc (1.5M); GIL mechanics and CPU-bound concurrency impact (2M).
\end{markingrubric}

\noindent\textbf{Question 17 [5 Marks]:} Implement a class \texttt{PolynomialRegressionPipeline} in pure Python and NumPy that takes coordinates $(x, y)$, fits an $n$-degree polynomial using normal equations, and evaluates predictions.
\begin{solutionbox}
\begin{lstlisting}
import numpy as np

class PolynomialRegressionPipeline:
    def __init__(self, degree: int):
        self.degree = degree
        self.weights = None

    def fit(self, x: np.ndarray, y: np.ndarray):
        # Construct Vandermonde matrix X
        X = np.vander(x, N=self.degree + 1, increasing=True)
        # Normal equations: (X^T X) w = X^T y
        self.weights = np.linalg.solve(X.T @ X, X.T @ y)

    def predict(self, x: np.ndarray) -> np.ndarray:
        X = np.vander(x, N=self.degree + 1, increasing=True)
        return X @ self.weights
\end{lstlisting}
\end{solutionbox}
\begin{markingrubric}
\textbf{Marking Scheme:} Vandermonde matrix formulation (2M); Normal equation solve via \texttt{np.linalg.solve} (2M); Prediction matrix multiplication (1M).
\end{markingrubric}

\noindent\textbf{Question 18 [5 Marks]:} Compare list comprehensions, generator expressions, and map/filter functions in terms of execution latency and memory allocation.
\begin{solutionbox}
$\bullet$ \textbf{List Comprehension:} Eagerly allocates entire list in heap memory. Optimized in C loop; fast for small collections ($\mathcal{O}(n)$ memory).\\
$\bullet$ \textbf{Generator Expression:} Lazily computes elements on demand via iterator protocol ($\mathcal{O}(1)$ memory).\\
$\bullet$ \textbf{Map/Filter:} Returns an iterator in Python 3. Slightly slower than comprehension due to Python function call overhead unless using built-in C functions.
\end{solutionbox}
\begin{markingrubric}
\textbf{Marking Scheme:} Eager vs lazy memory allocation comparison (2M); Complexity analysis (2M); Performance trade-offs (1M).
\end{markingrubric}

\noindent\textbf{Question 19 [5 Marks]:} Write a Python script using Pandas that loads a CSV with missing values, imputes missing numerical columns with their median, creates a normalized column, and outputs summary statistics.
\begin{solutionbox}
\begin{lstlisting}
import pandas as pd

def clean_and_normalize(filepath: str) -> pd.DataFrame:
    df = pd.read_csv(filepath)
    num_cols = df.select_dtypes(include=['float64', 'int64']).columns
    for col in num_cols:
        df[col] = df[col].fillna(df[col].median())
        min_v, max_v = df[col].min(), df[col].max()
        if max_v > min_v:
            df[f"{col}_norm"] = (df[col] - min_v) / (max_v - min_v)
    return df
\end{lstlisting}
\end{solutionbox}
\begin{markingrubric}
\textbf{Marking Scheme:} Reading CSV and selecting numeric columns (1M); Median imputation (2M); Min-max normalization (2M).
\end{markingrubric}

\noindent\textbf{Question 20 [5 Marks]:} Explain method overloading and method overriding in Python. How does Python support overloading without multiple method signatures?
\begin{solutionbox}
$\bullet$ \textbf{Overriding:} A subclass provides a specific implementation of a method defined in its superclass using identical names.\\
$\bullet$ \textbf{Overloading:} Python does NOT support traditional C++/Java method signature overloading because functions are first-class objects (last definition overwrites earlier ones). Overloading is achieved via default arguments (\texttt{arg=None}), variable arguments (\texttt{*args, **kwargs}), or \texttt{functools.singledispatch}.
\end{solutionbox}
\begin{markingrubric}
\textbf{Marking Scheme:} Clear definition of overriding (2M); Explanation of why Python lacks signature-based overloading and workarounds (3M).
\end{markingrubric}

\subsection*{Part C: In-Depth Engineering Systems \& Pipeline Problems (30 Marks)}

\noindent\textbf{Question 21 [10 Marks]:} Design and implement a complete Object-Oriented Financial Transaction Ledger System supporting atomic transfers, audit logs, and custom exceptions.
\begin{solutionbox}
\begin{lstlisting}
class LedgerException(Exception): pass
class BalanceError(LedgerException): pass

class Account:
    def __init__(self, owner: str, initial_balance: float):
        self.owner = owner
        self.balance = float(initial_balance)

    def deposit(self, amount: float):
        if amount <= 0: raise ValueError("Deposit must be positive")
        self.balance += amount

    def withdraw(self, amount: float):
        if amount > self.balance:
            raise BalanceError(f"Insufficient funds: Balance={self.balance}, Request={amount}")
        self.balance -= amount

class TransactionEngine:
    def __init__(self):
        self.audit_log = []

    def transfer(self, source: Account, destination: Account, amount: float):
        if amount <= 0: raise ValueError("Amount must be positive")
        # Atomic transfer execution
        source.withdraw(amount)
        try:
            destination.deposit(amount)
        except Exception as err:
            # Rollback
            source.deposit(amount)
            raise LedgerException(f"Transfer failed: {err}")
        self.audit_log.append(f"Transferred {amount} from {source.owner} to {destination.owner}")
\end{lstlisting}
\end{solutionbox}
\begin{markingrubric}
\textbf{Marking Scheme:} Custom exception classes (2M); Account encapsulation and balance validation (3M); TransactionEngine with atomic rollback logic (3M); Audit logging (2M).
\end{markingrubric}

\noindent\textbf{Question 22 [10 Marks]:} Implement Singular Value Decomposition (SVD) based dimensionality reduction and image compression using NumPy.
\begin{solutionbox}
\begin{lstlisting}
import numpy as np

def compress_matrix(A: np.ndarray, k: int):
    '''
    Compresses 2D matrix A using top k singular values.
    A = U @ diag(Sigma) @ Vt
    '''
    U, s, Vt = np.linalg.svd(A, full_matrices=False)
    # Truncate to top k singular vectors
    U_k = U[:, :k]
    s_k = s[:k]
    Vt_k = Vt[:k, :]
    
    A_compressed = U_k @ np.diag(s_k) @ Vt_k
    original_size = A.size
    compressed_size = U_k.size + s_k.size + Vt_k.size
    compression_ratio = compressed_size / original_size
    return A_compressed, compression_ratio
\end{lstlisting}
\textbf{Mathematical Explanation:} Eckart-Young-Mirsky Theorem guarantees that truncated SVD provides the optimal rank-$k$ approximation minimizing Frobenius norm $\|A - A_k\|_F$.
\end{solutionbox}
\begin{markingrubric}
\textbf{Marking Scheme:} SVD decomposition call via \texttt{np.linalg.svd} (3M); Truncation of $U, \Sigma, V^T$ to $k$ components (3M); Matrix reconstruction product (2M); Compression ratio derivation (2M).
\end{markingrubric}

\noindent\textbf{Question 23 [10 Marks]:} Build an end-to-end automated ETL and reporting pipeline using Pandas and Matplotlib for an IoT temperature sensor dataset.
\begin{solutionbox}
\begin{lstlisting}
import pandas as pd
import matplotlib.pyplot as plt

def execute_iot_pipeline(csv_file: str, report_plot: str):
    # 1. Ingestion & Cleaning
    df = pd.read_csv(csv_file, parse_dates=['timestamp'])
    df = df.dropna(subset=['temperature'])
    
    # 2. Resampling & Feature Extraction
    df.set_index('timestamp', inplace=True)
    hourly = df['temperature'].resample('1H').agg(['mean', 'max', 'min']).dropna()
    
    # 3. Multi-panel Visualization
    fig, (ax1, ax2) = plt.subplots(2, 1, figsize=(10, 8), dpi=300, sharex=True)
    ax1.plot(hourly.index, hourly['mean'], color='#0F2043', label='Hourly Mean Temp')
    ax1.fill_between(hourly.index, hourly['min'], hourly['max'], color='#EA580C', alpha=0.2, label='Min-Max Range')
    ax1.set_ylabel('Temperature (C)')
    ax1.legend()
    ax1.grid(True, linestyle='--')
    
    ax2.hist(df['temperature'], bins=30, color='#0F2043', edgecolor='white')
    ax2.set_xlabel('Temperature Distribution')
    ax2.set_ylabel('Observation Count')
    ax2.grid(True, linestyle='--')
    
    plt.tight_layout()
    plt.savefig(report_plot)
    plt.close()
    return hourly
\end{lstlisting}
\end{solutionbox}
\begin{markingrubric}
\textbf{Marking Scheme:} Timestamp parsing and missing value imputation (2.5M); Temporal resampling with aggregation (2.5M); Dual-panel Matplotlib dashboard generation (3M); Code style, comments, and PEP 8 standards (2M).
\end{markingrubric}

\newpage
\section{End-Term Comprehensive Examination --- Set B (Units 1 to 6)}
\noindent\textbf{Course Code:} CSR101 \hfill \textbf{Max Marks:} 70\\
\textbf{Course Title:} Python Programming (AI \& Data Engg.) \hfill \textbf{Time Allowed:} 3 Hours\\
\vspace{4mm}
\hrule
\vspace{4mm}

\subsection*{Part A: Objective Conceptual Questions (15 Marks)}
\begin{enumerate}[label=\textbf{\arabic*.}]
    \item What is the purpose of Python virtual environments (\texttt{venv})? \hfill \textbf{[Ans: Isolated dependency management per project]}
    \item What is printed by \texttt{print("a" * 3)}? \hfill \textbf{[Ans: \texttt{'aaa'}]}
    \item Which regex flag enables case-insensitive pattern matching? \hfill \textbf{[Ans: \texttt{re.IGNORECASE} or \texttt{re.I}]}
    \item What is the time complexity of deleting an element from the beginning of a Python list? \hfill \textbf{[Ans: $\mathcal{O}(n)$ due to shifting elements]}
    \item In Python, what is the base class for all built-in exceptions? \hfill \textbf{[Ans: \texttt{BaseException}]}
    \item What does the \texttt{zip()} function return? \hfill \textbf{[Ans: An iterator of tuples]}
    \item Which NumPy function flattens an array returning a contiguous copy? \hfill \textbf{[Ans: \texttt{flatten()} (unlike \texttt{ravel()} which returns a view)]}
    \item What is the shape of \texttt{np.array([[1, 2], [3, 4]])}? \hfill \textbf{[Ans: \texttt{(2, 2)}]}
    \item What operator is used for matrix multiplication in Python 3.5+? \hfill \textbf{[Ans: \texttt{@}]}
    \item In Pandas, how do you drop duplicate rows based on column \texttt{'id'}? \hfill \textbf{[Ans: \texttt{df.drop\_duplicates(subset=['id'])]}}
    \item What is the purpose of \texttt{plt.tight\_layout()} in Matplotlib? \hfill \textbf{[Ans: Automatically adjusts subplot params to prevent overlap]}
    \item What does \texttt{open("data.txt", "a")} do if the file does not exist? \hfill \textbf{[Ans: Creates a new empty file]}
    \item Which keyword rebinds a variable in the nearest enclosing non-global scope? \hfill \textbf{[Ans: \texttt{nonlocal}]}
    \item What is returned by \texttt{filter(None, [0, 1, False, 2, ''])}? \hfill \textbf{[Ans: Iterator yielding truthy values: \texttt{1, 2}]}
    \item What does \texttt{pd.merge(..., how='inner')} compute? \hfill \textbf{[Ans: Intersection of records matching key values]}
\end{enumerate}

\subsection*{Part B: Analytical \& Coding Problems (25 Marks)}

\noindent\textbf{Question 16 [5 Marks]:} Explain how Python executes list comprehensions internally. Compare list comprehensions with Python \texttt{for} loops and explain why comprehensions execute faster.
\begin{solutionbox}
List comprehensions are compiled into specialized bytecode leveraging the \texttt{LIST\_APPEND} instruction inside the C-level evaluation loop. In contrast, standard Python \texttt{for} loops incur repeated name lookup overhead for \texttt{list.append} attribute retrieval and multiple function call evaluation stack frames on each iteration.
\end{solutionbox}
\begin{markingrubric}
\textbf{Marking Scheme:} Bytecode instruction \texttt{LIST\_APPEND} (2.5M); Elimination of method lookup and call frame overhead (2.5M).
\end{markingrubric}

\noindent\textbf{Question 17 [5 Marks]:} Write a Python decorator \texttt{@retry(max\_attempts=3, delay=1.0)} that automatically retries a failing network function upon encountering an exception.
\begin{solutionbox}
\begin{lstlisting}
import time
from functools import wraps

def retry(max_attempts=3, delay=1.0):
    def decorator(func):
        @wraps(func)
        def wrapper(*args, **kwargs):
            attempts = 0
            while attempts < max_attempts:
                try:
                    return func(*args, **kwargs)
                except Exception as err:
                    attempts += 1
                    if attempts >= max_attempts:
                        raise err
                    time.sleep(delay)
        return wrapper
    return decorator
\end{lstlisting}
\end{solutionbox}
\begin{markingrubric}
\textbf{Marking Scheme:} Three-layer nested closure architecture for parameterized decorator (3M); Exception catching and retry sleep loop (2M).
\end{markingrubric}

\noindent\textbf{Question 18 [5 Marks]:} Implement a function \texttt{matrix\_power(A, p)} using NumPy and binary exponentiation ($\mathcal{O}(\log p)$ matrix multiplications).
\begin{solutionbox}
\begin{lstlisting}
import numpy as np

def matrix_power(A: np.ndarray, p: int) -> np.ndarray:
    n = A.shape[0]
    result = np.eye(n)
    base = A.copy()
    while p > 0:
        if p % 2 == 1:
            result = result @ base
        base = base @ base
        p //= 2
    return result
\end{lstlisting}
\end{solutionbox}
\begin{markingrubric}
\textbf{Marking Scheme:} Binary exponentiation logic (3M); Matrix multiplication with \texttt{@} (2M).
\end{markingrubric}

\noindent\textbf{Question 19 [5 Marks]:} Describe the split-apply-combine strategy in Pandas. Give a worked example computing the mean and 95th percentile latency grouped by API route.
\begin{solutionbox}
\begin{lstlisting}
import pandas as pd
import numpy as np

# Sample data
df = pd.DataFrame({
    'route': ['/login', '/api/v1/predict', '/login', '/api/v1/predict'],
    'latency': [45.2, 120.5, 52.1, 115.8]
})

res = df.groupby('route')['latency'].agg(
    mean_latency='mean',
    p95_latency=lambda x: np.percentile(x, 95)
).reset_index()
\end{lstlisting}
\end{solutionbox}
\begin{markingrubric}
\textbf{Marking Scheme:} Explanation of split-apply-combine steps (2M); GroupBy code with aggregation and lambda percentile (3M).
\end{markingrubric}

\noindent\textbf{Question 20 [5 Marks]:} Explain the difference between \texttt{\_\_str\_\_} and \texttt{\_\_repr\_\_}. Provide a class demonstrating both methods with distinct representations.
\begin{solutionbox}
$\bullet$ \texttt{\_\_repr\_\_}: Goal is to be unambiguous, intended for developers and debugging (e.g., \texttt{eval(repr(obj)) == obj}).\\
$\bullet$ \texttt{\_\_str\_\_}: Goal is to be readable, intended for end users (\texttt{print(obj)}).
\begin{lstlisting}
class User:
    def __init__(self, username, user_id):
        self.username, self.user_id = username, user_id

    def __repr__(self):
        return f"User('{self.username}', {self.user_id})"

    def __str__(self):
        return f"User: {self.username} (ID: #{self.user_id})"
\end{lstlisting}
\end{solutionbox}
\begin{markingrubric}
\textbf{Marking Scheme:} Distinguishing audience and purpose of each method (3M); Code implementation (2M).
\end{markingrubric}

\subsection*{Part C: In-Depth Engineering Systems \& Pipeline Problems (30 Marks)}

\noindent\textbf{Question 21 [10 Marks]:} Design and implement a Thread-Safe In-Memory Key-Value Store with TTL (Time-To-Live) expiration and an eviction worker thread in Python.
\begin{solutionbox}
\begin{lstlisting}
import time
import threading

class TTLKeyStore:
    def __init__(self):
        self._store = {}
        self._lock = threading.Lock()

    def set(self, key: str, value, ttl_seconds: float):
        with self._lock:
            expiry = time.time() + ttl_seconds
            self._store[key] = (value, expiry)

    def get(self, key: str):
        with self._lock:
            if key not in self._store:
                return None
            val, expiry = self._store[key]
            if time.time() > expiry:
                del self._store[key]
                return None
            return val

    def cleanup_expired(self):
        with self._lock:
            now = time.time()
            expired_keys = [k for k, (v, exp) in self._store.items() if now > exp]
            for k in expired_keys:
                del self._store[k]
            return len(expired_keys)
\end{lstlisting}
\end{solutionbox}
\begin{markingrubric}
\textbf{Marking Scheme:} Thread synchronization via \texttt{threading.Lock} (3M); Store implementation with timestamp metadata (3M); Lazy expiration on get (2M); Batch cleanup implementation (2M).
\end{markingrubric}

\noindent\textbf{Question 22 [10 Marks]:} Implement the Gram-Schmidt Orthogonalization algorithm in NumPy to calculate the QR decomposition ($A = Q R$) of a full-rank matrix.
\begin{solutionbox}
\begin{lstlisting}
import numpy as np

def gram_schmidt_qr(A: np.ndarray):
    m, n = A.shape
    Q = np.zeros((m, n))
    R = np.zeros((n, n))
    
    for j in range(n):
        v = A[:, j].astype(float)
        for i in range(j):
            R[i, j] = np.dot(Q[:, i], A[:, j])
            v = v - R[i, j] * Q[:, i]
        R[j, j] = np.linalg.norm(v)
        if R[j, j] < 1e-12:
            raise ValueError("Matrix is rank deficient.")
        Q[:, j] = v / R[j, j]
        
    return Q, R
\end{lstlisting}
\textbf{Orthogonality Check:} $Q^T Q = I_n$ and $A = Q R$.
\end{solutionbox}
\begin{markingrubric}
\textbf{Marking Scheme:} Mathematical formulation of projection vector subtraction (4M); Normalization and $Q, R$ construction (4M); Verification of orthogonality (2M).
\end{markingrubric}

\noindent\textbf{Question 23 [10 Marks]:} Build an anomaly detection pipeline on multivariate data using Pandas and Matplotlib. Calculate the Mahalanobis distance from scratch using NumPy.
\begin{solutionbox}
\begin{lstlisting}
import numpy as np
import pandas as pd
import matplotlib.pyplot as plt

def detect_mahalanobis_anomalies(df: pd.DataFrame, threshold: float = 3.0):
    X = df.values
    mu = np.mean(X, axis=0)
    X_diff = X - mu
    cov = np.cov(X, rowvar=False)
    inv_cov = np.linalg.inv(cov)
    
    # Compute Mahalanobis distance for each sample: d = sqrt((x-mu)^T S^-1 (x-mu))
    # Vectorized calculation:
    left = X_diff @ inv_cov
    dists = np.sqrt(np.sum(left * X_diff, axis=1))
    
    df['mahalanobis'] = dists
    df['is_anomaly'] = dists > threshold
    
    # Visualization
    fig, ax = plt.subplots(figsize=(8, 4), dpi=300)
    ax.scatter(np.arange(len(dists)), dists, c=df['is_anomaly'].map({True: '#DC3545', False: '#0F2043'}), alpha=0.7)
    ax.axhline(threshold, color='#EA580C', linestyle='--', label=f'Threshold ({threshold})')
    ax.set_ylabel('Mahalanobis Distance')
    ax.set_xlabel('Sample Index')
    ax.legend()
    plt.tight_layout()
    plt.savefig('anomalies.pdf')
    plt.close()
    return df
\end{lstlisting}
\end{solutionbox}
\begin{markingrubric}
\textbf{Marking Scheme:} Covariance computation and matrix inversion (3M); Vectorized Mahalanobis calculation (3M); Thresholding and Pandas annotation (2M); Visual scatter plot generation (2M).
\end{markingrubric}


\end{document}
